% ============================================================
% APUNTES DE ESTUDIO — DERECHOS PERSONALÍSIMOS
% Documento autónomo generado según método Cornell
% ============================================================

\documentclass[12pt,oneside]{book}

% ---------------------------------------------------------------
% METADATOS DEL DOCUMENTO
% ---------------------------------------------------------------
\newcommand{\universidad}{Universidad de Buenos Aires}
\newcommand{\facultad}{Facultad de Derecho}
\newcommand{\carrera}{}
\newcommand{\materia}{Derechos Personalísimos}
\newcommand{\titulo}{Introducción a la Parte General}
\newcommand{\autor}{Isaac Edgar Camacho Ocampo}
\newcommand{\anio}{2024-2025}

% ---------------------------------------------------------------
% PAQUETES
% ---------------------------------------------------------------
\usepackage[utf8]{inputenc}
\usepackage[T1]{fontenc}
\usepackage[spanish,es-nodecimaldot]{babel}

\usepackage[
    top=2.5cm,
    bottom=2.5cm,
    left=3cm,
    right=2.5cm,
    headheight=15pt
]{geometry}

\usepackage{amsmath}
\usepackage{amssymb}
\usepackage{mathtools}

\usepackage{graphicx}
\usepackage{float}
\usepackage{caption}
\usepackage{subcaption}

\usepackage{booktabs}
\usepackage{array}
\usepackage{longtable}
\usepackage{tabularx}

\usepackage{enumitem}

\usepackage{xcolor}
\usepackage[most]{tcolorbox}

\usepackage{fancyhdr}
\usepackage{ragged2e}

% ---------------------------------------------------------------
% RUTAS DE IMÁGENES
% ---------------------------------------------------------------
\graphicspath{
    {./img/}
    {./graficos/}
    {./figuras/}
}

% ---------------------------------------------------------------
% HIPERVÍNCULOS Y METADATOS PDF
% ---------------------------------------------------------------
\usepackage[
    colorlinks=true,
    linkcolor=blue!50!black,
    citecolor=green!40!black,
    urlcolor=blue!60!black,
    pdftitle={\titulo},
    pdfauthor={\autor},
    pdfsubject={Apuntes universitarios},
    bookmarks=true
]{hyperref}

% ---------------------------------------------------------------
% CONFIGURACIÓN GENERAL
% ---------------------------------------------------------------
\setcounter{tocdepth}{2}

\setlength{\parindent}{0pt}
\setlength{\parskip}{0.5em}

\setlist[itemize]{
    topsep=0.3em,
    itemsep=0.2em
}
\setlist[enumerate]{
    topsep=0.3em,
    itemsep=0.2em
}

\clubpenalty=10000
\widowpenalty=10000

% ---------------------------------------------------------------
% COMANDOS MATEMÁTICOS
% ---------------------------------------------------------------
\newcommand{\abs}[1]{\left|#1\right|}
\newcommand{\norm}[1]{\left\lVert#1\right\rVert}
\newcommand{\vect}[1]{\mathbf{#1}}
\newcommand{\deriv}[2]{\frac{d#1}{d#2}}
\newcommand{\pderiv}[2]{\frac{\partial#1}{\partial#2}}

% ---------------------------------------------------------------
% ENCABEZADOS Y PIES DE PÁGINA
% ---------------------------------------------------------------
\pagestyle{fancy}
\fancyhf{}

\fancyhead[L]{\nouppercase{\leftmark}}
\fancyhead[R]{\materia}

\fancyfoot[C]{\thepage}

\renewcommand{\headrulewidth}{0.4pt}
\renewcommand{\footrulewidth}{0pt}

\fancypagestyle{plain}{
    \fancyhf{}
    \fancyfoot[C]{\thepage}
    \renewcommand{\headrulewidth}{0pt}
}

% ---------------------------------------------------------------
% CAJAS ACADÉMICAS
% ---------------------------------------------------------------
\newtcolorbox{definition}[1]{
    enhanced,
    breakable,
    title={Definición: #1},
    fonttitle=\bfseries,
    colback=blue!3,
    colframe=blue!50!black,
    boxrule=0.8pt,
    arc=2mm
}

\newtcolorbox{important}{
    enhanced,
    breakable,
    title={Importante},
    fonttitle=\bfseries,
    colback=yellow!8,
    colframe=orange!70!black,
    boxrule=0.8pt,
    arc=2mm
}

\newtcolorbox{example}[1]{
    enhanced,
    breakable,
    title={Ejemplo: #1},
    fonttitle=\bfseries,
    colback=green!4,
    colframe=green!40!black,
    boxrule=0.8pt,
    arc=2mm
}

\newtcolorbox{formula}[1]{
    enhanced,
    breakable,
    title={Fórmula / Regla: #1},
    fonttitle=\bfseries,
    colback=purple!4,
    colframe=purple!50!black,
    boxrule=0.8pt,
    arc=2mm
}

\newtcolorbox{exercise}[1]{
    enhanced,
    breakable,
    title={Ejercicio: #1},
    fonttitle=\bfseries,
    colback=gray!5,
    colframe=gray!60!black,
    boxrule=0.8pt,
    arc=2mm
}

\newtcolorbox{note}{
    enhanced,
    breakable,
    title={Nota},
    fonttitle=\bfseries,
    colback=white,
    colframe=gray!50,
    boxrule=0.6pt,
    arc=2mm
}

% ================================================================
\begin{document}

% ---------------------------------------------------------------
% PORTADA
% ---------------------------------------------------------------
\begin{titlepage}

\thispagestyle{empty}

\begin{center}

\vspace*{1cm}

{\large \textsc{\universidad}}

\vspace{0.2cm}

{\large \textsc{\facultad}}

\vspace{3cm}

{\Huge \textbf{\materia}}

\vspace{1cm}

{\LARGE \titulo}

\vspace{0.8cm}

\textit{Comprender a fondo estos institutos es el primer paso para ejercer la abogacía con rigor y responsabilidad.}

\vspace{1.2cm}

\rule{0.70\textwidth}{0.5pt}

\vspace{0.5cm}

{\large Apuntes de estudio}

\vspace{0.5cm}

\rule{0.70\textwidth}{0.5pt}

\vfill

{\large \autor}

\vspace{0.3cm}

{\large \anio}

\end{center}

\vfill

\begin{flushleft}
\textit{Este es un modesto aporte a los estudiantes de ``\materia'' no pretende sustituir el material oficial de la materia.}
\end{flushleft}

\end{titlepage}

% ---------------------------------------------------------------
% ÍNDICE
% ---------------------------------------------------------------
\tableofcontents

% ================================================================
% CAPÍTULO ÚNICO
% ================================================================
\chapter[Derechos Personalísimos]{Derechos Personalísimos: Introducción a la Parte General}

\section{Introducción}

Los derechos personalísimos son derechos subjetivos fundamentales que protegen la dignidad humana en sus manifestaciones físicas y espirituales. Estos derechos reconocen y garantizan a cada persona el respeto a su vida, integridad corporal, libertad, honor, intimidad e identidad.

En el ejercicio profesional de la abogacía, la comprensión de estos derechos resulta necesaria para:

\begin{itemize}
    \item Defender a clientes cuyos derechos personalísimos han sido vulnerados.
    \item Asesorar sobre el consentimiento informado en procedimientos médicos.
    \item Proteger la intimidad y la privacidad de datos en el entorno digital.
    \item Litigar casos de difamación, invasión de la privacidad e injurias.
    \item Orientar cuestiones relacionadas con la bioética y los experimentos humanos.
    \item Comprender la protección constitucional de la dignidad humana.
\end{itemize}

Estos derechos adquieren especial relevancia en el contexto de la tecnología digital avanzada, donde los datos personales, las imágenes y la información íntima pueden ser capturados, almacenados y difundidos sin consentimiento. El marco legal que se desarrolla a continuación permite proteger tanto los derechos propios como los ajenos frente a estos desafíos contemporáneos.

\section{Mapa de contenidos de la unidad}

Los derechos personalísimos forman un sistema integrado en el que la dignidad humana constituye el principio rector. El desarrollo de la unidad comienza con los conceptos generales y los antecedentes, continúa con las características y la clasificación, prosigue con la regulación en el código vigente y, finalmente, profundiza en dos dimensiones clave: la integridad física (que protege el cuerpo y la vida) y la integridad espiritual (que protege el honor, la imagen, la intimidad y la identidad). Cada bloque se construye sobre el anterior, de forma análoga a las capas que envuelven un núcleo central de protección de la persona humana.

\begin{note}
El presente documento corresponde a la primera de tres clases dedicadas a esta unidad temática. Los bloques de \textbf{Integridad física} e \textbf{Integridad espiritual} del cuadro siguiente serán desarrollados en clases posteriores; aquí se cubren los bloques de \textbf{Fundamentos y conceptos generales}, \textbf{Características y clasificación}, \textbf{Regulación en el Código Civil y Comercial} y una primera aproximación al \textbf{Derecho a la libertad}.
\end{note}

\begin{longtable}{p{0.30\textwidth} p{0.62\textwidth}}
\toprule
\textbf{Tema} & \textbf{Contenido y pregunta orientadora} \\
\midrule
\endfirsthead
\toprule
\textbf{Tema} & \textbf{Contenido y pregunta orientadora} \\
\midrule
\endhead
\multicolumn{2}{l}{\textit{Fundamentos y conceptos generales}} \\
\addlinespace
1. Noción de Derechos Personalísimos &
Definición, derechos subjetivos, interioridad, manifestaciones físicas y espirituales. \newline
\textit{¿Qué son los derechos personalísimos y por qué son tan importantes para la persona humana?} \\
\addlinespace
2. Importancia y contexto actual &
Circunstancias contemporáneas, biotecnologías, datos, privacidad digital. \newline
\textit{¿Cuáles son las nuevas amenazas y desafíos del siglo XXI respecto a los derechos personalísimos?} \\
\addlinespace
3. Antecedentes en Argentina &
Código de Vélez, evolución legislativa, leyes especiales, reformas. \newline
\textit{¿Cómo ha evolucionado la protección de estos derechos en la legislación argentina?} \\
\addlinespace
\multicolumn{2}{l}{\textit{Características y clasificación}} \\
\addlinespace
4. Caracteres de los Derechos Personalísimos &
Subjetivos, innatos, vitalicios, necesarios, absolutos, inherentes, intransmisibles, irrenunciables. \newline
\textit{¿Cuáles son las características fundamentales que definen a los derechos personalísimos?} \\
\addlinespace
5. Clasificación tripartita &
Integridad física, Libertad, Integridad espiritual. \newline
\textit{¿Cómo se clasifican los derechos personalísimos según el tipo de bien protegido?} \\
\addlinespace
6. Dignidad humana como fundamento &
Dimensión ontológica, dimensión de autonomía, concepto integrador. \newline
\textit{¿Cuál es el concepto central que articula la protección de todos los derechos personalísimos?} \\
\addlinespace
\multicolumn{2}{l}{\textit{Regulación en el Código Civil y Comercial}} \\
\addlinespace
7. Disposiciones generales &
Artículos 51-62, medios de protección, prevención y reparación. \newline
\textit{¿Cuáles son los mecanismos de protección que establece el código civil para estos derechos?} \\
\addlinespace
8. Disponibilidad relativa &
Artículo 55, consentimiento, restricciones, revocabilidad. \newline
\textit{¿En qué casos se pueden disponer los derechos personalísimos y bajo qué condiciones?} \\
\addlinespace
9. Derecho a la libertad &
Artículo 19 CN, actos privados, protección civil y constitucional. \newline
\textit{¿Cómo protege el ordenamiento jurídico el derecho fundamental a la libertad?} \\
\addlinespace
\multicolumn{2}{l}{\textit{Integridad física} (clase posterior)} \\
\addlinespace
10. Vida e integridad corporal &
Artículos 54, 56, 57, actos peligrosos, disposición del cuerpo, embriones. \newline
\textit{¿Qué límites establece el código respecto a los actos que pueden comprometer la integridad física?} \\
\addlinespace
11. Consentimiento informado &
Artículo 59, requisitos, información, revocabilidad, directivas anticipadas. \newline
\textit{¿Cuáles son los requisitos para un consentimiento válido en materia de procedimientos médicos?} \\
\addlinespace
12. Trasplantes e investigación humana &
Ley 27.447, principio de totalidad, solidaridad, requisitos, protecciones. \newline
\textit{¿Cómo se equilibra la disposición corporal con la protección de la dignidad en trasplantes e investigación?} \\
\addlinespace
\multicolumn{2}{l}{\textit{Integridad espiritual} (clase posterior)} \\
\addlinespace
13. Derecho al honor &
Honor subjetivo y objetivo, fama, reputación, calumnias e injurias civiles. \newline
\textit{¿Cómo se protege el honor y la reputación de una persona en el ámbito civil?} \\
\addlinespace
14. Derecho a la imagen &
Captura y reproducción, consentimiento, excepciones, personas fallecidas. \newline
\textit{¿Bajo qué circunstancias se puede captar, reproducir y difundir la imagen de una persona?} \\
\addlinespace
15. Derecho a la intimidad &
Artículo 1770, injerencias arbitrarias, privacidad, datos personales, habeas data. \newline
\textit{¿Cómo se protege la intimidad y la privacidad ante intromisiones injustificadas?} \\
\addlinespace
16. Derecho a la identidad &
Identidad biológica y social, filiación, orígenes, características dinámicas. \newline
\textit{¿Cuál es la importancia del derecho a conocer la propia identidad y origen?} \\
\bottomrule
\end{longtable}

\section{Noción de derechos personalísimos}

Los derechos personalísimos corresponden a la persona humana. Son derechos que pertenecen a todo ser humano por el solo hecho de ser tal.

\begin{definition}{Derechos personalísimos (Tobías)}
Los derechos personalísimos son derechos subjetivos que reconocen y garantizan a la persona humana el goce en su propia entidad, en su interioridad, en todas sus manifestaciones físicas y espirituales.
\end{definition}

Estos derechos no deben confundirse con los atributos de la personalidad, distinción que tampoco realiza el Código Civil y Comercial. Los atributos de la personalidad son cualidades esenciales que permiten que la persona despliegue su personalidad y entable relaciones jurídicas. Los derechos personalísimos, en cambio, son verdaderos derechos subjetivos que protegen dimensiones fundamentales de la persona.

Estos derechos protegen dos dimensiones fundamentales:

\begin{enumerate}
    \item La \textbf{dimensión de la integridad física}: la vida y la integridad corporal.
    \item La \textbf{dimensión espiritual}: la libertad, el honor, la imagen, la intimidad y la identidad de la persona.
\end{enumerate}

\section{Importancia y contexto actual}

A lo largo de la historia siempre ha existido una protección de los valores esenciales de la persona: la vida, la integridad física, la libertad, la imagen y el honor. Sin embargo, especialmente en la última mitad del siglo XX, se dieron circunstancias que motivaron a los juristas y al legislador a sancionar normas y a adoptar tratados internacionales orientados a proteger estas dimensiones de la persona humana.

Han surgido nuevas formas de amenaza que reclaman una respuesta jurídica. En el entorno de las biotecnologías, hoy es posible intervenir sobre el cuerpo humano de una manera mucho más poderosa: el conocimiento de la genética y de la biología ha dado lugar a enormes beneficios, pero también a nuevas formas de intervención que pueden no ser respetuosas de la integridad física.

\begin{note}
A lo largo de la historia se han realizado experimentos sobre personas humanas sin su consentimiento, aprovechando la condición de vulnerabilidad de determinados grupos, lo que ha dado lugar a fuertes reacciones éticas y jurídicas.
\end{note}

Las nuevas tecnologías inciden también en la transmisión de la vida humana, en los trasplantes de órganos y en las intervenciones biomédicas y biotecnológicas en general. Respecto de la investigación en seres humanos, la genética permite hoy intervenciones de gran alcance, como ocurrió en China en el año 2018, cuando se obtuvieron seres genéticamente modificados. Estas posibilidades tecnológicas son enormemente promisorias para el bienestar y el avance de las personas, pero presentan también aspectos problemáticos que requieren una reflexión cuidadosa.

\subsection{Datos personales y privacidad digital}

Otro ámbito relevante es el de la recolección, el uso, el tratamiento, la transmisión y el almacenamiento de los datos personales: datos vinculados con la identidad, con las características personales, con el consumo o con la salud. La cantidad de datos que hoy se almacena es enorme, fenómeno conocido como \textit{big data}: la tendencia a generar grandes colecciones de datos de las personas para elaborar perfiles, realizar seguimientos y dirigir ofertas publicitarias.

Los datos personales permiten tomar decisiones más racionales y conocer mejor a las personas, lo cual puede traer beneficios, pero también nuevas posibilidades de discriminación y de ataques a las personas en función de sus datos. Cuando se trata de datos sensibles existen nuevas posibilidades de manipulación; así, en determinadas situaciones políticas, las campañas se han dirigido a públicos específicos para persuadirlos e influir en su voto.

\begin{important}
El desarrollo de la informática ha transformado radicalmente la capacidad de procesamiento de datos personales —antes limitada a fichas escritas a mano— lo que exige nuevas formas de protección de la intimidad, la identidad y el honor de las personas. La posibilidad actual de captar la voz y la imagen de manera mucho más poderosa hace que estos elementos, una vez registrados, se conviertan también en datos personales sujetos a este mismo régimen de protección.
\end{important}

\section{Antecedentes en Argentina}

Entre los autores que en la Argentina han desarrollado el estudio de esta materia se encuentran Salvat, Cifuentes, Rivera, Navarro Floria y Tobías, entre otros. Cabe señalar que este tema desborda el contenido propio de la parte general del derecho civil y podría constituir por sí mismo una materia completa.

% TODO: contenido incompleto o ambiguo en las notas originales — el segundo autor mencionado aparece en el apunte como "el profesor Salto Cifuentes", lo que podría corresponder a dos autores distintos (Salvat y Cifuentes) mal unidos en la transcripción; se conserva la grafía más plausible sin poder confirmarlo con el material disponible.

\subsection{El Código de Vélez}

El Código de Vélez no contenía un tratamiento metódico y orgánico de los derechos personalísimos. No obstante, los autores reconocen la presencia del tema en algunos de sus artículos:

\begin{longtable}{p{0.30\textwidth} p{0.62\textwidth}}
\toprule
\textbf{Disposición} & \textbf{Contenido} \\
\midrule
\endfirsthead
\toprule
\textbf{Disposición} & \textbf{Contenido} \\
\midrule
\endhead
Artículos 428 y 2.126 & Hablan de derechos inherentes a la persona. \\
\addlinespace
Artículo 1.075 & Establece que toda afectación de un derecho que se confunda con la existencia de la persona puede ser materia de un delito. \\
\bottomrule
\end{longtable}

También en la nota al artículo 2.312, el Código de Vélez hacía referencia a ciertos derechos que tienen origen en la existencia del individuo mismo, como la libertad, el honor, el cuerpo y la patria potestad. Se trata de una mención presente en el código, aunque no sistemática ni ordenada.

\subsection{Evolución legislativa en Argentina}

\begin{longtable}{p{0.30\textwidth} p{0.62\textwidth}}
\toprule
\textbf{Norma} & \textbf{Contenido} \\
\midrule
\endfirsthead
\toprule
\textbf{Norma} & \textbf{Contenido} \\
\midrule
\endhead
Ley 11.723 (1933) — Ley de Propiedad Intelectual & Norma que continúa vigente. En sus artículos 31 a 35 regula la fotografía y el derecho a la imagen, incluyendo el consentimiento y los casos en que puede prescindirse de él. \\
\addlinespace
Ley 21.173 — incorporó el artículo 1.071 bis al código anterior & Protegía el derecho a la intimidad. Este artículo sería recogido posteriormente por el nuevo Código Civil y Comercial en su artículo 1.770. \\
% TODO: contenido incompleto o ambiguo en las notas originales — el apunte indica el año 1905 para esta incorporación, lo cual resulta inconsistente con la numeración de la ley citada; se conserva el dato tal como figura en el material original.
\addlinespace
Ley 24.193 (1992) — Ley de Trasplantes & Reemplazada en 2018 por la Ley 27.447. \\
\addlinespace
Ley 25.326 (2000) — Ley de Protección de Datos Personales & Aún vigente, con proyectos de reforma. \\
\addlinespace
Ley 26.529 (2009) — Ley de Derechos del Paciente & Reformada en 2012 por la Ley 26.742. \\
\bottomrule
\end{longtable}

Estas normas reflejan las dimensiones de los derechos personalísimos vinculadas con la integridad física, la integridad espiritual y la libertad.

\subsection{La Constitución Nacional de 1994}

La reforma constitucional de 1994 incorpora, en el artículo 75 inciso 22, diversos tratados de derechos humanos con jerarquía constitucional que consagran derechos correspondientes a la personalidad.

Se produce así una convergencia entre el derecho público y el derecho privado. Aunque los tratados de derechos humanos incorporados a la Constitución corresponden al derecho público, tienen resonancias directas en el derecho privado, en tanto corresponden a la persona humana, que se despliega fundamentalmente en el campo civil.

\begin{note}
La protección de estos derechos tiene también faceta penal: la protección del derecho al honor puede dar lugar al delito de calumnias e injurias, además de a la protección civil. Ambas protecciones suelen coexistir. El presente desarrollo se limita a las dimensiones civiles del asunto y a su incorporación al Código Civil y Comercial.
\end{note}

\section{Caracteres de los derechos personalísimos}

Según Tobías, los derechos personalísimos presentan los siguientes caracteres:

\begin{itemize}
    \item \textbf{Son derechos subjetivos privados:} corresponden a la esfera de las personas, dimensión donde operan el derecho privado y el derecho civil.
    \item \textbf{Son derechos innatos y originarios:} no son derechos concedidos por el legislador, sino reconocidos. Un ser humano, por el solo hecho de ser tal, tiene derecho a la vida, al honor, a la imagen y a la identidad.
    \item \textbf{Son vitalicios y necesarios:} acompañan a la persona durante toda su vida; una persona humana siempre debe tenerlos.
    \item \textbf{Son absolutos (o de exclusión):} tienen un sujeto pasivo universal. Frente a cualquier persona existe la obligación de respetarlos, sin que sea necesario haber celebrado un contrato ni haber tenido relación previa alguna.
    \item \textbf{Son inherentes a la persona:} no pueden separarse de la persona; forman parte de su ser mismo.
    \item \textbf{Son intransmisibles:} no es posible transmitir completamente, por ejemplo, el derecho a la integridad física, ni delegar en otro la toma de decisiones sobre estos derechos.
    \item \textbf{Son relativamente indisponibles:} en principio no se pueden disponer, salvo mediante disposiciones relativas, de carácter restrictivo y siempre con posibilidad de revocación.
    \item \textbf{Son irrenunciables:} no es posible renunciar a estos derechos. La manifestación de ceder o renunciar, por ejemplo, al derecho a la imagen, carece de validez.
\end{itemize}

\begin{important}
Estos caracteres protegen dimensiones fundamentales de la persona y, en última instancia, tienen como bien jurídico común la dignidad humana.
\end{important}

\section{Clasificación tripartita}

Se adopta aquí la clasificación tripartita habitualmente utilizada por el profesor Tobías, sin perjuicio de que otros autores propongan clasificaciones distintas.

\begin{table}[H]
\centering
\begin{tabular}{p{0.28\textwidth} p{0.28\textwidth} p{0.28\textwidth}}
\toprule
\textbf{Integridad física} & \textbf{Libertad} & \textbf{Integridad espiritual} \\
\midrule
Vida \newline Integridad corporal \newline Salud &
Concepto que comprende aspectos tanto corporales como espirituales &
Honor \newline Intimidad \newline Imagen \newline Identidad \\
\bottomrule
\end{tabular}
\end{table}

En torno a esta clasificación se organiza el desarrollo de la unidad. En la presente clase se aborda una primera aproximación a la libertad y a las disposiciones generales del Código Civil y Comercial; las clases siguientes estarán dedicadas a la integridad física y a la integridad espiritual.

\section{Regulación en el Código Civil y Comercial}

Metodológicamente, el Código Civil y Comercial incorpora los derechos personalísimos en el Capítulo 3, Título Primero del Libro Primero. El Libro Primero —Parte General— se abre con el Título Primero, dedicado a la persona humana; dentro de este título, luego de regular el comienzo de la existencia de la persona y la capacidad de las personas, aparecen los derechos personalísimos en los artículos 51 a 62.

\begin{important}
La incorporación sistemática de los derechos personalísimos constituye una de las grandes novedades del Código Civil y Comercial, en la medida en que el Código de Vélez no traía un tratamiento sistemático de estos derechos. Esta incorporación otorga una mayor presencia e importancia a la dignidad humana en el tratamiento del Código Civil y Comercial.
\end{important}

\subsection{Disposiciones generales: artículo 51}

El capítulo se abre con el artículo 51, referido a la inviolabilidad de la persona humana:

\begin{formula}{Artículo 51 CCyC}
``La persona humana es inviolable y en cualquier circunstancia tiene derecho al reconocimiento y respeto de su dignidad.''
\end{formula}

El artículo no menciona explícitamente el derecho a la vida, como sí lo hacían proyectos anteriores, pero dicha protección puede reconocerse detrás de la noción de inviolabilidad de la persona: no es posible afectar su persona, ni su cuerpo, ni su integridad física, ni su integridad espiritual. La expresión ``reconocimiento'' señala que la dignidad no constituye una concesión del ordenamiento jurídico, sino algo preexistente que el ordenamiento se limita a reconocer.

\subsection{Medios de protección: prevención y reparación}

El artículo 52 establece que la enumeración de medios de protección no es taxativa, sin perjuicio de que exista una sistematización que distingue entre:

\begin{table}[H]
\centering
\begin{tabular}{p{0.30\textwidth} p{0.55\textwidth}}
\toprule
\textbf{Prevención} & \textbf{Reparación} \\
\midrule
Consiste en evitar que ocurra el hecho que menoscabaría la dignidad; se trata de impedir que se produzca el acto dañino. &
Si el hecho ocurriere, se trata de volver al estado anterior y, si ello no fuera posible, de indemnizar los daños producidos. \\
\bottomrule
\end{tabular}
\end{table}

El Código incorpora, además, nuevas formas de reparación, entre ellas la posibilidad de publicar la sentencia en la que se ha discutido la afectación de un derecho personalísimo. Esta forma de reparación, presente ya en artículos anteriores, se incorpora también en la parte de daños del Libro Tercero como posibilidad específica cuando se encuentra afectado el honor de la persona.

\subsection{Disponibilidad relativa: artículo 55}

El artículo 55 dispone que los derechos personalísimos no son disponibles, salvo mediante el consentimiento de la persona, consentimiento que se encuentra sujeto a una serie de límites.

\begin{formula}{Artículo 55 CCyC — Disponibilidad relativa}
El consentimiento no puede ser contrario a las buenas costumbres. El consentimiento no se presume; es de interpretación restrictiva y, ante la duda, se interpreta que no ha sido otorgado. Es libremente revocable en cualquier momento.
\end{formula}

\begin{important}
El consentimiento sobre los derechos personalísimos debe reunir las siguientes condiciones:
\begin{itemize}
    \item No puede ser contrario a la ley, la moral o las buenas costumbres.
    \item Debe ser expresado; no puede presumirse a partir de la mera falta de oposición. Puede manifestarse por escrito, en forma verbal o mediante un gesto comprensible, pero siempre debe existir una manifestación real.
    \item Es de interpretación restrictiva: ante la duda, se interpreta que no hubo consentimiento. Así, el consentimiento otorgado para el uso de la propia imagen en una campaña de bien común no se extiende, por interpretación restrictiva, a una campaña con fines comerciales.
    \item Es libremente revocable: la persona puede dar marcha atrás y revocar el consentimiento oportunamente otorgado.
\end{itemize}
\end{important}

\section{Derecho a la libertad}

El valor de la libertad constituye uno de los derechos personalísimos más importantes. Su fuente se ubica generalmente en el artículo 19 de la Constitución Nacional.

\begin{formula}{Artículo 19 de la Constitución Nacional}
``Las acciones privadas de los hombres que de ningún modo ofendan a la moral pública, ni perjudiquen a un tercero, están exentas de la autoridad de los magistrados. Ningún habitante de la Nación será obligado a hacer lo que la ley no manda, ni privado de lo que ella no prohíbe.''
\end{formula}

Este constituye un principio de libertad fundamental en el ordenamiento jurídico argentino. Desde la perspectiva civil, la protección de la libertad se vincula sobre todo con los supuestos en que una persona sufre violencia para celebrar un acto; en tal caso, la persona afectada tiene derecho a reclamar los daños y perjuicios correspondientes.

También existe protección de la libertad en el campo de los derechos constitucionales, a través del habeas corpus, de la acción de amparo y de la acción de inconstitucionalidad, instrumentos que permiten conocer la situación de una persona, su paradero y su situación respecto de su integridad y de su libertad.

\begin{note}
Al igual que en el resto de los derechos personalísimos, las formas de protección de la libertad se articulan a través del artículo 52, mediante la prevención y la reparación.
\end{note}

% ================================================================
% CUADRO CORNELL
% ================================================================
\section{Cuadro Cornell de repaso}

El siguiente cuadro permite repasar los conceptos centrales desarrollados en esta unidad mediante la técnica de recuperación activa propia del método Cornell.

\begin{longtable}{
    >{\raggedright\arraybackslash}p{0.30\textwidth}
    >{\raggedright\arraybackslash}p{0.60\textwidth}
}
\toprule
\textbf{Preguntas / palabras clave} & \textbf{Notas / respuestas} \\
\midrule
\endfirsthead
\toprule
\textbf{Preguntas / palabras clave} & \textbf{Notas / respuestas} \\
\midrule
\endhead

\textbf{¿Qué son los derechos personalísimos?} &
Son derechos subjetivos que reconocen y garantizan a la persona humana el goce en su propia entidad, en su interioridad, en todas sus manifestaciones físicas y espirituales. Protegen dimensiones fundamentales del ser humano inherentes a su naturaleza y dignidad. \\
\addlinespace

\textbf{¿Cuál es la diferencia entre derechos personalísimos y atributos de la personalidad?} &
Los atributos de la personalidad son cualidades esenciales que permiten que la persona despliegue su personalidad y entable relaciones jurídicas (como el nombre, el domicilio o la capacidad). Los derechos personalísimos son verdaderos derechos subjetivos que protegen dimensiones fundamentales de la persona (como la vida, el honor o la integridad física). \\
\addlinespace

\textbf{¿Por qué es importante estudiar los derechos personalísimos hoy?} &
Las circunstancias contemporáneas —biotecnologías, big data, privacidad digital— han generado nuevas amenazas a estos derechos que requieren una respuesta jurídica actualizada, orientada a proteger a las personas frente a formas de intromisión y vulneración que no existían en épocas anteriores. \\
\addlinespace

\textbf{¿Cuáles son los caracteres principales de los derechos personalísimos?} &
Son: subjetivos privados; innatos y originarios (reconocidos, no concedidos); vitalicios y necesarios; absolutos (sujeto pasivo universal); inherentes a la persona; intransmisibles; relativamente indisponibles; e irrenunciables. Todos estos caracteres buscan proteger la dignidad humana. \\
\addlinespace

\textbf{¿Cómo se clasifican los derechos personalísimos?} &
En forma tripartita: integridad física (vida, integridad corporal, salud); libertad (concepto que comprende aspectos corporales y espirituales); e integridad espiritual (honor, intimidad, imagen, identidad). \\
\addlinespace

\textbf{¿Dónde se regulan los derechos personalísimos en el Código Civil y Comercial?} &
En el Capítulo 3, Título Primero, Libro Primero (artículos 51 a 62). El artículo 51 abre el capítulo con el principio de inviolabilidad de la persona humana y su derecho al reconocimiento y respeto de su dignidad. \\
\addlinespace

\textbf{¿Cuáles son los medios de protección de los derechos personalísimos?} &
El artículo 52 establece: la prevención, que busca evitar que ocurra el acto dañino, y la reparación, que busca volver al estado anterior o indemnizar. También se prevé la publicación de la sentencia como forma específica de reparación en casos de daño al honor, la intimidad o la identidad. \\
\addlinespace

\textbf{¿Qué significa que los derechos personalísimos son relativamente indisponibles?} &
No pueden ser dispuestos en principio, salvo mediante consentimiento que: no sea contrario a las buenas costumbres; no se presuma (debe ser expresado); sea de interpretación restrictiva; y sea libremente revocable. \\
\addlinespace

\textbf{¿Cómo se protege la libertad en el derecho civil?} &
A través de acciones civiles cuando existe violencia o coacción para celebrar un acto (reclamo de daños y perjuicios). También existen protecciones constitucionales, como el habeas corpus, el amparo y la acción de inconstitucionalidad, más específicas respecto de la libertad física. \\
\addlinespace

\midrule
\multicolumn{2}{p{0.92\textwidth}}{
\textbf{Resumen:} Los derechos personalísimos constituyen un sistema integrado de protección de la dignidad humana en todas sus dimensiones. Fundados en el reconocimiento de que cada persona posee una dignidad inherente, protegen la vida, la integridad física y espiritual, la libertad, el honor, la intimidad y la identidad de todo ser humano. Aunque tienen raíces profundas en la historia del derecho, adquieren una importancia renovada en el contexto contemporáneo, en el que las biotecnologías, la digitalización y el almacenamiento masivo de datos han generado nuevas amenazas a estas esferas fundamentales de la persona. El Código Civil y Comercial supone un avance significativo al sistematizar estos derechos en su Parte General, reconociendo su carácter innato, vitalicio, necesario, absoluto, inherente, intransmisible e irrenunciable. Su protección se articula mediante dos vías principales: la prevención y la reparación.
} \\
\bottomrule
\end{longtable}

\end{document}
